\begin{frame}{Ratios de gestión}
\color{gestion}\textbf{Eficiencia en el uso de activos y administración del ciclo operativo}\color{black}

\vspace{0.7em}
Permiten revisar si las ventas se convierten en efectivo oportunamente y si los inventarios permanecen demasiado tiempo inmovilizados.
\end{frame}

\begin{frame}{Rotación de cuentas por cobrar}
\[
\text{Rotación CxC}=\frac{\text{Ventas netas a crédito}}{\text{Promedio de cuentas por cobrar}}
\]
\textbf{Datos necesarios:} ventas a crédito, cuentas por cobrar iniciales y finales.
\end{frame}

\begin{frame}{Periodo promedio de cobro}
\[
\text{Periodo promedio de cobro}=\frac{360}{\text{Rotación de cuentas por cobrar}}
\]
Se usan 360 días como criterio didáctico contable. También puede usarse 365 días, siempre que se mantenga consistencia.
\end{frame}

\begin{frame}{Caso práctico: cobranza}
\casoeducativo
\small
Ventas a crédito: \soles{1,200,000.00}. CxC inicial: \soles{180,000.00}. CxC final: \soles{220,000.00}.
\[
\text{Promedio}=\frac{180,000.00+220,000.00}{2}=200,000.00
\]
\[
\text{Rotación}=\frac{1,200,000.00}{200,000.00}=6.00
\quad
\text{Periodo}=\frac{360}{6.00}=60.00\text{ días}
\]
\textbf{Decisión:} comparar 60 días con la política de crédito.
\end{frame}

\begin{frame}{Rotación de inventarios}
\[
\text{Rotación de inventarios}=\frac{\text{Costo de ventas}}{\text{Inventario promedio}}
\]
Evalúa cuántas veces se renueva el inventario durante el periodo.
\end{frame}

\begin{frame}{Caso práctico: inventarios}
\casoeducativo
\small
Costo de ventas: \soles{900,000.00}. Inventario inicial: \soles{140,000.00}. Inventario final: \soles{160,000.00}.
\[
\text{Inventario promedio}=\frac{140,000.00+160,000.00}{2}=150,000.00
\]
\[
\text{Rotación}=\frac{900,000.00}{150,000.00}=6.00
\quad
\text{Permanencia}=\frac{360}{6.00}=60.00\text{ días}
\]
\textbf{Interpretación:} el inventario permanece cerca de 60 días en almacén.
\end{frame}
